\documentclass[10pt,a4paper]{article}

\usepackage[utf8]{inputenc}
\usepackage[margin=18mm]{geometry}
\usepackage{times}
\usepackage{microtype}
\usepackage{titlesec}
\usepackage{enumitem}
\usepackage{tabularx}
\usepackage{hyperref}
\usepackage{xcolor}

% Color definitions
\definecolor{linkcolor}{RGB}{21, 62, 103}
\definecolor{rulecolor}{RGB}{34, 34, 34}

% Hyperref setup
\hypersetup{
    colorlinks=true,
    linkcolor=linkcolor,
    urlcolor=linkcolor,
    citecolor=linkcolor,
    pdftitle={Prince Vats - Academic Curriculum Vitae},
    pdfauthor={Prince Vats}
}

% Section formatting
\titleformat{\section}
  {\large\bfseries\scshape\raggedright}
  {}{0em}
  {}[\titlerule]

\titlespacing*{\section}{0pt}{10pt}{5pt}

% Subsection formatting
\titleformat{\subsection}
  {\normalsize\bfseries\raggedright}
  {}{0em}
  {}

\titlespacing*{\subsection}{0pt}{6pt}{3pt}

% Custom list environment
\setlist[itemize]{leftmargin=14pt, itemsep=2pt, topsep=2pt, parsep=0pt}

\setlength{\parindent}{0pt}
\setlength{\parskip}{3pt}

\begin{document}

% ==================== HEADER ====================
\begin{center}
    {\LARGE \textbf{\textsc{Prince Vats}}}\\[3pt]
    {\large \textbf{Ph.D. Researcher in Environmental Engineering}}\\[2pt]
    {\normalsize \textit{Department of Civil \& Environmental Engineering $\bullet$ Indian Institute of Technology Gandhinagar}}\\[4pt]
    \small
    Email: \href{mailto:princevatsiitrpr@gmail.com}{princevatsiitrpr@gmail.com} $\;\bullet\;$
    Phone: +91-9027902245 $\;\bullet\;$
    Location: Haridwar / Gandhinagar, India\\[2pt]
    Website: \href{https://prince-vats.github.io}{prince-vats.github.io} $\;\bullet\;$
    \href{https://scholar.google.com/citations?user=SPHQZR8AAAAJ&hl=en}{Google Scholar} $\;\bullet\;$
    \href{https://www.linkedin.com/in/prince-vats-iitrpr/}{LinkedIn} $\;\bullet\;$
    \href{https://github.com/prince-vats}{GitHub}
\end{center}

\vspace{2pt}

% ==================== RESEARCH FOCUS ====================
\section*{Research Focus \& Core Interests}
\textbf{Aerosol Science \& Particle Dynamics:} Particle number size distribution (PNSD) reconstruction from light-scattering measurements, aerosol microphysics, atmospheric transformation, and human respiratory exposure.\\
\textbf{Environmental IoT Networks:} Design, multi-pollutant sensing ($\text{PM}_1, \text{PM}_{2.5}, \text{PM}_{10}, \text{CO}_2, \text{CO}, \text{TVOCs}$), automated telemetry, and machine-learning field calibration against reference monitors.\\
\textbf{Spatio-Temporal Deep Learning:} Advanced neural network architectures (CNNs, GNNs, LSTMs, GRUs) for high-resolution dynamic forecasting, spatial air pollution interpolation, and pattern recognition.\\
\textbf{Health Risk Modeling:} Deterministic and stochastic environmental exposure analysis, indoor air quality mitigation via phytoremediation, and epidemiological health outcome assessments.

% ==================== EDUCATION ====================
\section*{Education}

\textbf{Indian Institute of Technology Gandhinagar} \hfill \textbf{2024 -- Present (Ongoing)}\\
\textit{Ph.D. Candidate, Environmental Engineering} \hfill \textit{Gandhinagar, Gujarat, India}
\begin{itemize}
    \item \textbf{Doctoral Research:} AI-driven spatio-temporal analysis and high-resolution prediction of urban air quality using dense IoT-based environmental monitoring networks.
    \item Development and field deployment of low-cost Environmental Quality Monitoring Systems (EQMS) for high-frequency urban monitoring.
    \item Machine learning pipelines for field sensor calibration, data fusion, and exposure modeling.
    \item Deep learning models (GNN, LSTM/GRU, CNN-BiLSTM) for short-term and multi-scale $\text{PM}_{2.5}$ forecasting.
\end{itemize}

\vspace{3pt}
\textbf{Indian Institute of Technology Ropar} \hfill \textbf{2022 -- 2024}\\
\textit{Master of Technology (M.Tech), Environmental Engineering} \hfill \textit{Rupnagar, Punjab, India}
\begin{itemize}
    \item \textbf{M.Tech Thesis:} Investigated indoor air pollution ($\text{PM}_{10}, \text{PM}_{2.5}, \text{PM}_1$, and black carbon) across classrooms, hostels, and offices; $\text{PM}_{2.5}$ ($59.6\,\mu\text{g/m}^3$) and $\text{PM}_{10}$ ($71.12\,\mu\text{g/m}^3$) substantially exceeded CPCB permissible limits.
    \item \textbf{Phytoremediation Strategy:} Evaluated Ajwain (\textit{Trachyspermum ammi}, 39.3\% $\text{PM}_{2.5}$ reduction) and Money plant (\textit{Epipremnum aureum}, 38.7\% $\text{PM}_{10}$ reduction), establishing practical room-scale botanical mitigation recommendations.
    \item \textbf{Respiratory Health Assessment:} Conducted clinical spirometry on 303 students, identifying 35.97\% restrictive patterns, 34.65\% obstructive patterns, and 74.59\% asthma prevalence linked to particulate exposure.
\end{itemize}

\vspace{3pt}
\textbf{G. B. Pant University of Agriculture and Technology (GBPUAT)} \hfill \textbf{2018 -- 2022}\\
\textit{Bachelor of Technology (B.Tech), Civil Engineering} \hfill \textit{Pantnagar, Uttarakhand, India}

% ==================== HONORS & AWARDS ====================
\section*{Honors \& National Competitive Examinations}
\begin{itemize}
    \item \textbf{GATE 2025:} \textbf{All India Rank (AIR) 29} (99.56 percentile) nationwide in India's premier Graduate Aptitude Test in Engineering.
    \item \textbf{GATE 2026:} \textbf{All India Rank (AIR) 32} (99.56 percentile) nationwide.
    \item \textbf{Best Paper Award:} 7th National Water Seminar (7NWS), National Institute of Hydrology (NIH), Roorkee (August 2023).
\end{itemize}

% ==================== PUBLICATIONS ====================
\section*{Peer-Reviewed Publications}

\subsection*{Journal Articles}
\begin{enumerate}[leftmargin=18pt, label={[\arabic*]}]
    \item Dhada, I., \textbf{Vats, P.}, and Samal, S. K. (2025). Potential Health Risk Estimation Through Deterministic and Stochastic Approaches for the Air Pollutants Found Inside Bottling Beverage Industries of Himachal Pradesh in India. \textit{Discover Environment}, Springer, March 2025.\\
    DOI: \href{https://doi.org/10.1007/s44274-025-00221-x}{10.1007/s44274-025-00221-x}
\end{enumerate}

\subsection*{Refereed Book Chapters \& Conference Proceedings}
\begin{enumerate}[leftmargin=18pt, label={[\arabic*]}, resume]
    \item Dhada, I., \textbf{Vats, P.}, and Singh, B. (2025). Probable Health Risk of Pollutants and Noise Found Inside Bottling Beverage Plant of Himachal Pradesh, India. In: Shiva Nagendra, S.M., Ramamurthy, P.C., Vardoulakis, S., Tantrakarnapa, K., Ball, R.J. (eds) \textit{Select Proceedings of the 8th Indian International Conference on Air Quality Management (IICAQM 2023)}. \textit{Lecture Notes in Civil Engineering}, vol 582, pp. 219--238. Springer, Singapore. \textit{(Oral Presentation at IICAQM 2023)}.\\
    DOI: \href{https://doi.org/10.1007/978-981-96-2359-4_18}{10.1007/978-981-96-2359-4\_18}
\end{enumerate}

% ==================== CONFERENCES ====================
\section*{Conference Presentations \& Scholarly Contributions}

\subsection*{Peer-Reviewed Oral Presentations}
\begin{enumerate}[leftmargin=18pt, label={[\arabic*]}]
    \item \textbf{Vats, P.}, Thakur, A. K., and Patel, S. (2025, Dec). ``Development and Deployment of a Low-Cost IoT Based Environmental Quality Monitoring System (EQMS) with Spatio-Temporal Predictive Modeling'', \textit{14th Asian Aerosol Conference (AAC 2025)}, Indian Aerosol Science and Technology Association (IASTA), Mumbai, India. \textit{(Presented by P. Vats)}.
    \item Dhada, I., \textbf{Vats, P.}, and Singh, B. (2023, Dec). ``Probable Health Risk of Pollutants and Noise Found Inside Bottling Beverage Plant of Himachal Pradesh, India'', \textit{8th Indian International Conference on Air Quality Management (IICAQM 2023)}, IISc Bengaluru, India. \textit{(Presented by I. Dhada; Published in Springer LNCE, 2025)}.
\end{enumerate}

\subsection*{Refereed Poster Presentations}
\begin{enumerate}[leftmargin=18pt, label={[\arabic*]}, resume]
    \item \textbf{Vats, P.}, Nagar, C., Sen, S., and Patel, S. (2026, Apr). ``Low-Cost Environmental Quality Monitoring System for High-Resolution Spatiotemporal Urban Air Quality Monitoring and Multi-Scale Forecasting'', \textit{Industry Partnership Conclave (IPC 2026)}, IIT Gandhinagar, India.
    \item \textbf{Vats, P.}, Nagar, C., Sen, S., and Patel, S. (2026, Apr). ``Environmental Quality Monitoring System Using Low-Cost Sensors for Urban Air Quality Monitoring and Multi-Scale Forecasting'', \textit{KPCSD Poster Competition on Sustainability}, Kiran C. Patel Centre for Sustainable Development, IIT Gandhinagar, India.
    \item Singh, A. P., \textbf{Vats, P.}, and Jain, M. K. (2025, Dec). ``ML Model Evaluation for AQI Prediction in IGP Using Air Quality and Meteorological Features'', \textit{14th Asian Aerosol Conference (AAC 2025)}, Mumbai, India.
    \item \textbf{Vats, P.}, and Patel, S. (2025, Oct). ``Environmental Quality Monitoring System Using Low-Cost Sensors for Urban Air Quality Monitoring and Multi-Scale Forecasting'', \textit{Chemical Engineering Symposium (CES 2025)}, IIT Gandhinagar \& SVNIT, India.
    \item \textbf{Vats, P.}, Singh, A. P., and Patel, S. (2025, Aug). ``Evaluating CNN-LSTM and GNN-GRU Hybrid Models for Enhanced Air Quality Prediction and Dynamic Forecasting'', \textit{India Clean Air Summit (ICAS 2025)}, CSTEP, Bengaluru, India.
    \item Singh, A. P., \textbf{Vats, P.}, and Jain, M. K. (2025, Aug). ``ML Model Evaluation for AQI Prediction in Gangetic Plain Using Air Quality and Meteorological Features'', \textit{India Clean Air Summit (ICAS 2025)}, CSTEP, Bengaluru, India.
    \item \textbf{Two Paper Presentations} (2023, Aug). \textit{7th National Water Seminar (7NWS)}, National Institute of Hydrology (NIH), Roorkee, India. \textit{(Conferred Best Paper Award)}.
\end{enumerate}

% ==================== TEACHING ====================
\section*{Teaching \& Academic Mentorship}

\textbf{Fundamentals of Aerosol Science} --- \textit{Teaching Assistant} \hfill \textbf{AY 2025--2026, Semester II}\\
\textit{Indian Institute of Technology Gandhinagar $\bullet$ Advanced Undergraduate \& Postgraduate Level}
\begin{itemize}
    \item Assisted in instruction of core aerosol physics and atmospheric chemistry: particle mechanics, dynamics, size distribution functions (PNSD), and atmospheric transformation processes.
    \item Conducted problem-solving tutorial sessions, guided students on atmospheric datasets, and instructed on measurement instrumentation including Scanning Mobility Particle Sizers (SMPS) and Optical Particle Counters (OPC).
\end{itemize}

\vspace{3pt}
\textbf{ES 101: Engineering Graphics} --- \textit{Graduate Teaching Fellow} \hfill \textbf{AY 2025--2026 \& AY 2026--2027}\\
\textit{Indian Institute of Technology Gandhinagar $\bullet$ First-Year B.Tech ($\sim$50 Students per cohort)}
\begin{itemize}
    \item Taught classical engineering drafting: orthographic projections, isometric views, and sectioning of solids.
    \item Conducted weekly laboratory sessions instructing 3D parametric CAD modeling in \textbf{Autodesk Inventor} and \textbf{Autodesk Fusion 360}; evaluated design coursework and mentored student mechanical projects.
\end{itemize}

\vspace{3pt}
\textbf{Introduction to Writing} --- \textit{Teaching Assistant} \hfill \textbf{AY 2024--2025, Semester I}\\
\textit{Indian Institute of Technology Gandhinagar $\bullet$ First-Year B.Tech Foundation Course}
\begin{itemize}
    \item Mentored undergraduate students in academic and technical communication, structured argumentation, critical reading, and peer-review workflows; provided detailed editorial critique on technical writing.
\end{itemize}

% ==================== RESEARCH PROJECTS ====================
\section*{Research Projects \& Field Studies}

\textbf{Environmental Quality Monitoring System (EQMS) for Urban Air Quality} \hfill \textit{IIT Gandhinagar}
\begin{itemize}
    \item Developed an integrated, low-cost multi-pollutant sensing network measuring $\text{PM}_1, \text{PM}_{2.5}, \text{PM}_{10}, \text{CO}_2, \text{CO}$, TVOCs, and ambient meteorology (temperature, relative humidity, barometric pressure).
    \item Engineered end-to-end data pipeline streaming real-time measurements across a dense urban grid to a centralized database server.
    \item Implemented CNN-BiLSTM architectures for anomaly detection and multi-scale air quality forecasting across 24-hour, 7-day, and 1-month horizons to generate actionable environmental policy insights.
\end{itemize}

\vspace{2pt}
\textbf{Institutional Indoor Air Quality, Phytoremediation \& Spirometry Assessment} \hfill \textit{IIT Ropar}
\begin{itemize}
    \item Conducted longitudinal indoor air quality monitoring across classrooms, hostels, and offices, demonstrating high exposure risks.
    \item Benchmarked botanical mitigation efficacy of indoor plant species and administered clinical spirometry testing on 303 human subjects to establish correlations between particulate exposure and respiratory health.
\end{itemize}

\vspace{2pt}
\textbf{Occupational Exposure \& Health Risk Assessment in Beverage Bottling Plants} \hfill \textit{Himachal Pradesh}
\begin{itemize}
    \item Monitored TVOCs, $\text{PM}_{10}, \text{PM}_{2.5}, \text{PM}_1$, black carbon (BC), and industrial noise across 12 manufacturing facilities during monsoon.
    \item Formulated deterministic and stochastic health risk models, identifying non-carcinogenic hazard ratios (0.52--0.80) and carcinogenic BC risks (5.8E-04 to 6.8E-04) to guide industrial ventilation recommendations.
\end{itemize}

% ==================== ACADEMIC SERVICE ====================
\section*{Academic Service \& Event Organization}
\begin{itemize}
    \item \textbf{Organizing Team Member:} Faculty Development Program (FDP) for ASEAN and Global South, IIT Gandhinagar (September 2026).
    \item \textbf{Organizing Team Member:} Industry Partnership Conclave (IPC 2026), IIT Gandhinagar (April 2026).
    \item \textbf{Organizing Team Member:} Symposium on Promoting Aerosol Education and Research (PAER 2026), IIT Gandhinagar (February 2026).
    \item \textbf{Organizing Team Member:} Symposium on Advances in Tools and Technologies to Quantify Health Impacts of Air Pollution (APHES 2025), IIT Gandhinagar (December 2025).
\end{itemize}

% ==================== OPEN SOURCE TOOLS ====================
\section*{Scientific Software \& Open Research Tools}
\begin{itemize}
    \item \textbf{Atmospheric Gas Unit Converter:} Online scientific utility for converting ambient gaseous concentration units ($\text{ppb}$, $\mu\text{g/m}^3$) across user-defined temperature and atmospheric pressure regimes. (\href{https://prince-vats.github.io/unit-converter/}{prince-vats.github.io/unit-converter})
    \item \textbf{Sensor Calibration Tool for Atmospheric Measurements:} Interactive web-based data calibration framework designed for calibrating low-cost environmental sensors against reference-grade instrumentation. (\href{https://prince-vats.github.io/sensor-calibration/}{prince-vats.github.io/sensor-calibration})
\end{itemize}

% ==================== TECHNICAL TOOLKIT ====================
\section*{Technical \& Computational Toolkit}
\begin{tabularx}{\textwidth}{@{}p{0.26\textwidth}X@{}}
    \textbf{Programming \& Software:} & Python, C, C++, MATLAB, R, MS Office Suite \\
    \textbf{AI \& Deep Learning:} & CNN, Graph Neural Networks (GNN), LSTM, GRU, Transformers / Predictors, TensorFlow, PyTorch, Feature Engineering (PCA) \\
    \textbf{Spatial \& CAD Prototyping:} & ArcGIS, QGIS, Autodesk Inventor, Autodesk Fusion 360 \\
    \textbf{Specialized Domain Context:} & IoT Sensor Design \& Network Deployment, Environmental Data Calibration, Deterministic \& Stochastic Health Risk Assessments, Real-Time Air Quality Monitoring ($\text{PM}_{2.5}, \text{PM}_{10}, \text{TVOCs}$) \\
    \textbf{Languages:} & English (Fluent), Hindi (Fluent), Japanese (Basic)
\end{tabularx}

% ==================== WORKSHOPS ====================
\section*{Specialized Academic Workshops Attended}
\begin{itemize}
    \item \textbf{Smart Sensors and Data Analytics for Air Quality Management Workshop:} Indian Institute of Technology Bombay (March 2025).
    \item \textbf{AI4Water: Advanced Data-Driven and Artificial Intelligence Tools in Water Resources Workshop:} Indian Institute of Technology Roorkee (August 2023).
\end{itemize}

% ==================== LEADERSHIP ====================
\section*{Leadership \& Community Engagement}
\begin{itemize}
    \item \textbf{Vice President:} Green Club --- spearheaded environmental initiatives, clean-up drives, and campus sustainability projects.
    \item \textbf{Community Mentorship:} Taught underprivileged students through an NGO initiative.
    \item \textbf{Technical Seminars \& Hackathons:} Organized university technical seminars and actively participated in environmental hackathons.
    \item \textbf{Campus Sustainability:} Volunteered in community clean-up drives and environmental tree plantation initiatives.
\end{itemize}

\end{document}
